\section{Reuleaux Geometry and the Membrane Body}
\label{sec:geometry}

\subsection{Force Angle and shortfall}
Three circles of radius $s$ centred at the vertices of an equilateral triangle of
side $d$ produce
\[
\Delta\theta(d):=\frac{2\pi}{3}-\theta(d)=2\arcsin\Bigl(\frac{d}{2s}\Bigr)>0
\quad\text{for all }d>0.
\]
At the tetrahedral scale $d=s$ one has the critical landmark
\[
\Delta\theta_c=\frac{\pi}{3}.
\]

\subsection{Four-sphere membrane}
Lifting to four spheres of radius $s$ at the vertices of a regular tetrahedron of
side $s$ yields a closed membrane shell:
\begin{align*}
r_{\mathrm{inner}}&=\Bigl(1-\frac{\sqrt{6}}{4}\Bigr)s,\\
r_{\mathrm{outer}}&=\Bigl(1+\frac{\sqrt{6}}{4}\Bigr)s,\\
\Delta r_m&=\frac{\sqrt{6}}{2}s.
\end{align*}
The membrane is the region inside some but not all spheres---the three-dimensional
analogue of the two-but-not-three arc condition.

\subsection{Layered particle}
Overlap count partitions the particle into:
\begin{itemize}
\item \textbf{Heart} --- 4-of-4 overlap (maximal bonding; primary SNF locus);
\item \textbf{Skin} --- predominantly 3-of-4;
\item \textbf{Membrane} --- 1-of-4 / 2-of-4 (imperfect, leaky).
\end{itemize}
